\documentclass[11pt,a4paper]{article}
\usepackage[margin=2.5cm]{geometry}
\usepackage[utf8]{inputenc}
\usepackage[T1]{fontenc}
\usepackage[french]{babel}
\usepackage{enumitem}
\usepackage{hyperref}

\begin{document}
	\pagestyle{empty}
	
	\noindent
	Friedly WOLI \\
	Toulouse, 31100 \\
	0767339290 \\
	friedlysedjro@gmail.com
	
	\begin{flushright}
		Airbus \\
		Blagnac Wings Campus\\
		JR10393010
	\end{flushright}
	
	\bigskip
	
	\textbf{Objet : Candidature alternance – ingénieur(e) en méthodes numériques pour la mécanique des solides }
	
	\bigskip
	
Madame, Monsieur,\\

Actuellement étudiant en troisième année de mécanique énergétique à l'Université Toulouse III Paul Sabatier, je souhaite intégrer à partir de septembre 2026 un Master en modélisation et simulation en mécanique en alternance. Dans ce cadre, je vous adresse ma candidature pour l'offre d'apprentissage intitulée « Apprenti(e) ingénieur(e) en méthodes numériques avancées pour la mécanique des solides » au sein de Airbus sur le site de Blagnac.

Ma formation m’a permis d’acquérir de solides bases en mécanique des solides, résistance des matériaux, mécanique des fluides et modélisation mathématique des systèmes physiques. J’ai également développé des compétences en programmation scientifique, notamment en Python, à travers plusieurs projets de travaux dirigés consacrés à l’analyse de données expérimentales et à la visualisation de résultats sous forme de courbes et de cartes de contours. Ces projets m’ont permis de développer une approche rigoureuse de l’analyse numérique et de la modélisation de phénomènes physiques. Par ailleurs, la rédaction de rapports scientifiques en LaTeX fait partie intégrante de ces travaux, renforçant mes compétences en rédaction scientifique.

En parallèle de mes études, je suis membre de l’association TLS'e Racing, basée à l’ISAE-SUPAERO, qui développe un véhicule électrique pour la compétition Formula Student 2026. Dans ce cadre, j’ai participé au choix du moteur ainsi qu’à la définition de la configuration de la batterie. J’ai également réalisé des simulations structurelles avec Ansys afin d’évaluer la résistance du conteneur de batterie aux charges mécaniques et aux accélérations dans différentes directions, en respectant les contraintes du cahier des charges de la compétition.

L’opportunité de rejoindre le département de recherche et technologies d’Airbus représente pour moi une occasion unique de contribuer à l’amélioration des modèles numériques utilisés pour prédire le comportement mécanique des matériaux. Les thématiques liées à la simulation avancée, à la mécanique des solides et à l’amélioration des modèles numériques pour des applications aéronautiques correspondent pleinement à mes objectifs professionnels et à mon intérêt pour la simulation scientifique appliquée aux structures complexes.

Rigoureux, curieux et motivé par les problématiques de modélisation et de simulation mécanique, je souhaite mettre mes compétences au service de vos travaux de recherche tout en développant une expertise en méthodes numériques avancées et en modélisation des matériaux.

Je vous remercie par avance de l’attention portée à ma candidature et me tiens à votre disposition pour un entretien.

Je vous prie d’agréer, Madame, Monsieur, l’expression de mes salutations distinguées.
	\vspace{1cm}
	
	\begin{flushright}
		\textbf{Friedly WOLI}
	\end{flushright}
	
\end{document}